% 编译引擎与字体 -------------------------------------------------
\usepackage{fontspec}
\setmainfont{TeX Gyre Termes}
\setsansfont{TeX Gyre Heros}
\setmonofont{Noto Sans Mono CJK SC}
\setCJKmainfont{Noto Serif CJK SC}
\setCJKsansfont{Noto Sans CJK SC}
\setCJKmonofont{Noto Sans Mono CJK SC}
\setCJKfamilyfont{zhkai}{AR PL KaitiM GB}
\newcommand{\kaiti}{\CJKfamily{zhkai}}

% 页面与段落 -----------------------------------------------------
\usepackage[a4paper, top=2.6cm, bottom=2.6cm, left=2.8cm, right=2.8cm,
            headheight=15pt, headsep=0.7cm, footskip=1.0cm]{geometry}
\usepackage{setspace}
\onehalfspacing
\setlength{\parindent}{2em}
\setlength{\parskip}{0pt}
\raggedbottom

\ctexset{
  section = {
    format = \centering\zihao{3}\bfseries,
    beforeskip = 1.6ex plus .4ex minus .2ex,
    afterskip = 1.2ex plus .2ex
  },
  subsection = {
    format = \zihao{4}\bfseries,
    beforeskip = 1.3ex plus .3ex minus .2ex,
    afterskip = .8ex plus .2ex
  },
  subsubsection = {
    format = \zihao{-4}\bfseries,
    beforeskip = 1.0ex plus .2ex,
    afterskip = .5ex
  }
}

% 页眉页脚 -------------------------------------------------------
\usepackage{fancyhdr}
\pagestyle{fancy}
\fancyhf{}
\fancyhead[L]{\zihao{-5}\CourseName\ 课程论文}
\fancyhead[R]{\zihao{-5}\PaperShortTitle}
\fancyfoot[C]{\zihao{-5}\thepage}
\renewcommand{\headrulewidth}{0.4pt}
\renewcommand{\footrulewidth}{0pt}
\fancypagestyle{plain}{
  \fancyhf{}
  \fancyfoot[C]{\zihao{-5}\thepage}
  \renewcommand{\headrulewidth}{0pt}
}

% 数学、图表与列表 -----------------------------------------------
\usepackage{amsmath, amssymb, mathtools}
\usepackage{graphicx}
\graphicspath{{figures/}}
\usepackage{booktabs, longtable, tabularx, multirow, array}
\usepackage{float}
\usepackage{caption}
\usepackage{subcaption}
\captionsetup{font=small, labelsep=quad}
\captionsetup[figure]{name=图}
\captionsetup[table]{name=表}
\usepackage{enumitem}
\setlist{nosep, leftmargin=2em}

% 代码 -----------------------------------------------------------
\usepackage{xcolor}
\definecolor{codebg}{HTML}{F6F8FA}
\definecolor{codecomment}{HTML}{57606A}
\definecolor{codekeyword}{HTML}{8250DF}
\definecolor{codestring}{HTML}{0A3069}
\usepackage{listings}
\lstset{
  basicstyle=\small\ttfamily,
  backgroundcolor=\color{codebg},
  keywordstyle=\color{codekeyword}\bfseries,
  commentstyle=\color{codecomment},
  stringstyle=\color{codestring},
  numbers=left,
  numberstyle=\scriptsize\color{codecomment},
  frame=single,
  rulecolor=\color{black!15},
  breaklines=true,
  showstringspaces=false,
  tabsize=2,
  columns=fullflexible
}

% 参考文献：GB/T 7714—2015，顺序编码制 ---------------------------
\usepackage{csquotes}
\usepackage[
  backend=biber,
  style=gb7714-2015,
  sorting=none,
  maxbibnames=99,
  giveninits=true
]{biblatex}

% 超链接应放在大多数宏包之后加载 ---------------------------------
\usepackage{xurl}
\usepackage[
  colorlinks=true,
  linkcolor=black,
  citecolor=black,
  urlcolor=blue!55!black,
  bookmarksnumbered=true,
  pdfstartview=FitH
]{hyperref}
\urlstyle{same}
\usepackage[capitalise, noabbrev]{cleveref}
\crefname{figure}{图}{图}
\crefname{table}{表}{表}
\crefname{equation}{式}{式}
\crefname{section}{节}{节}

% 写作阶段提示 ---------------------------------------------------
\usepackage[most]{tcolorbox}
\newif\ifdraftmode
\draftmodetrue
\newtcolorbox{guidancebox}{
  colback=blue!3,
  colframe=blue!35!black,
  boxrule=.5pt,
  arc=1mm,
  left=2mm,
  right=2mm,
  top=1mm,
  bottom=1mm,
  fonttitle=\bfseries,
  title=写作提示
}
\newcommand{\writingguide}[1]{%
  \ifdraftmode
    \begin{guidancebox}\small #1\end{guidancebox}
  \fi
}
\newcommand{\placeholder}[1]{%
  \ifdraftmode\textcolor{blue!55!black}{【#1】}\fi
}

% PDF 元数据在读取 metadata.tex 后由正文设置，避免中文书签乱码。
\AtBeginDocument{%
  \hypersetup{
    pdftitle={\PaperShortTitle},
    pdfauthor={\AuthorName},
    pdfsubject={\TrackNameEN}
  }
}
